\scanpage{86}
% archive sheet number: 85
% block text | lines 1-2
Soient $\widehat{O},\widehat{A}$ les complétés de $O,A$, alors $\Gamma$ y opère de façon canonique. On a
% block display | line 3
\[
\widehat{O}=\bigoplus_i\widehat{O}_{m_i}\qquad\widehat{A}=\bigoplus_i\underbrace{\widehat{A}_{m_i}}_{B_i}
\]
% block text | lines 4-6
et $\Gamma_d$ \unsure{invarie} $B_0$, tandis que $\Gamma\bmod\unsure{\Gamma_i|\Gamma_d}$ permute les $B_i$ entre eux de façon simplement transitive. Il en résulte que
% block display | line 7
\[
H^1(\Gamma,\widehat{A}^*)\simeq H^1(\Gamma_d,\widehat{A}_m^*)
\]
% block text | lines 8-11
d'après un lemme bien connu (valable aussi en cohomologie non commutative). Il s'ensuit, puisque $f$ et $g$ définissent le $=$ élément de $H^1(\Gamma_d,\widehat{A}_m^*)$, \deleted{qu'il existe un $u\in A_m^*$} \inserted{[on regarde le diagramme commutatif}
% block text | left of diagram
\deleted{tel que}
% block diagram | lines 12-14
\diagram{\begin{tikzcd}[ampersand replacement=\&]
H^1(\Gamma,A^*)\arrow[r]\arrow[d]\&H^1(\Gamma_d,A_m^*)\arrow[d]\\
H^1(\Gamma,\widehat{A}^*)\arrow[r,"\sim"]\&H^1(\Gamma_d,\widehat{A}_m^*)
\end{tikzcd}}
% block text | lines 14-16
\inserted{]} qu'ils ont $=$ image dans $H^1(\Gamma,\widehat{A})$ donc qu'il existe un $u\in\widehat{A}^*$ tel qu'on ait
% block display | line 17
\[
g(\unsure{\gamma|\sigma})=u^{-1}f(\gamma)\gamma(u)
\]
% block text | lines 18-19 cancelled block
\deleted{Par \unsure{conséquent}, on peut, pour \unsure{$\forall$|toute} puissance $\unsure{\mathfrak r|\mathfrak m|\mathfrak b}^{N}$ du radical de $O$, trouver $v\in A$ tel que}
% block display | line 20 cancelled block
\[
\deleted{g(\gamma)\equiv u^{-1}f(\gamma)\gamma(u)\quad\bmod\mathfrak r^N A}
\]
% block text | line 21
Considérons l'application \inserted{$\unsure{O^\Gamma}$-}\deleted{\unsure{linéaire}}
% block display | line 22
\[
\varphi:A\longrightarrow A^\Gamma
\]
% block text | line 23
définie par
% block display | line 23
\[
\varphi(u)=\bigl(ug(\gamma)-f(\gamma)\gamma(u)\bigr)_{\gamma\in\Gamma}
\]
